\section{实现、数据与运行记录}\label{app:implementation}
\subsection{程序实现与操作}
用一个完整问题串联实现：选房价，输入“用整体质量和居住面积预测成交价，先用线性回归”。检索器查找这几个字段和训练方法说明；DeepSeek API 生成包含目标、特征、划分与算法的分析方案；用户确认后，执行器拟合预处理与回归，再保存验证预测。
\begin{center}\small
\begin{tabularx}{\textwidth}{lY}
\toprule
文件 & 作用\\
\midrule
\code{tablelab/catalog.py} & 载入数据与字段说明，检查文件指纹，提供目标和变量信息。\\
\code{tablelab/retrieval.py} & 建立字符 TF--IDF 索引，对当前问题返回相关说明。\\
\code{tablelab/cloud/deepseek_online.py} & 组织 DeepSeek 请求，校验返回的 JSON 分析参数。\\
\code{tablelab/cloud/deepseek_transport.py} & 通过 HTTPS 调用 DeepSeek API，解析响应与错误。\\
\code{tablelab/contracts.py} & 规定参数结构及允许的操作，检查字段、算法与数量。\\
\code{tablelab/engine.py} & 筛选、统计、划分、预处理、拟合、评估和预测。\\
\code{tablelab/plotting.py} & 用实际结果生成 PDF 图。\\
\bottomrule
\end{tabularx}
\end{center}
请求使用 JSON 输出格式，并附带本地方案的 JSON Schema；返回值还须通过本地参数校验。例如算法必须填已实现名称，数值筛选符号必须来自固定集合；对于“低价房”这样未给阈值的提问，仍需要澄清具体含义。格式通过检查后，也应核对语义是否符合原问题。

用户接着说“改成随机森林”，模型可以读取上一轮方案，将算法改为森林，并保留已经确定的目标和特征。重新执行时产生新的运行编号，重新拟合模型，保存新的验证预测。若说“只看 2000 年以后建成的房屋”，还需要修改筛选条件，并在筛选样本上重新划分和训练。

改变样本后，前后两组误差的评价对象也变了。界面应同时显示这次的筛选条件和训练、验证样本数，便于区分“算法变了”和“被预测的房屋变了”。用户可以修改这些条件，再检查新的计算结果。

先问 CollgCr 的平均成交价，核对参与计算的 150 条记录。随后用面积和质量建立线性回归，观察训练与验证数量、均值基线和误差；再改用随机森林，比较同一批验证记录。最后输入一套房屋的字段值，查看预测以及是否存在未见类别或超出训练范围的数值。

切换到 Adult 时，先比较教育组的高收入样本比例，再建立逻辑回归，解释一条概率怎样经过 0.5 阈值变成类别。NBA 可以比较后续赛季留出与姓名分组；学生任务可以比较考前信息与明确标注的同期估计问题。这些步骤是现场演示的操作安排。

每轮保存 \code{plan.json} 和 \code{evidence.json}，分别记录分析参数与检索片段；\code{provenance.json} 记录数据、代码和环境信息；统计值写入 \code{result.json} 与 CSV。训练任务还保存 \code{split.csv}、\code{validation_predictions.csv}、预处理参数、系数或置换重要性，以及 PDF 诊断图。

这些文件可以回答三类具体追问：这次选了哪些记录，程序怎样变换变量，预测值如何得到。\code{reproduce.py} 根据保存方案重算，\code{executor_source.py} 保留当时的执行代码。复现需匹配数据和依赖版本；单条预测依赖当前会话模型，服务重启后须先按训练方案重新拟合。

讲解计算脚本 \code{tools/teaching_calculations.py} 从上传数据重算社区均价和三条记录的直线，用保存的参数重建房价和 Adult 预测，并用数值差分检查逻辑回归梯度与 Hessian。十二个基准模型沿用已保存的结果，讲解计算的输出保存在单独目录。

\subsection{数据处理记录}
上传的 Adult、NBA、学生 ZIP 及其内部文件哈希与下载清单一致。单独上传的房价 CSV 与 \code{house-prices.zip} 中的 CSV 逐字节相同。哈希用于识别文件版本。房价副本的下载出处及其与 Kaggle 官方副本的一致性待核实。原 ZIP、解压原件、处理副本分开保存，处理副本新增稳定行号。

房价有 6,965 个空白单元格，全部 1,460 行均至少含一个空白；Adult 两个文件合计有 6,465 个缺失单元格；NBA 与学生表未发现空白。原始行全部保留，没有全表删缺失、自动去重或分位数截尾。训练时只对所选特征拟合插补规则。

房价 \code{MasVnrType} 有 864 个字面值 \code{None} 与 8 个空白，NBA 的 \code{college} 有 1,854 个字面值 \code{None}。处理时保留这些类别字符串。房价 \code{LotFrontage} 缺失 259 条；地下室质量的 37 个空白均对应地下室面积为零，壁炉质量的 690 个空白均对应无壁炉，车库类型的 81 个空白均对应车位数为零，泳池质量的 1,453 个空白均对应泳池面积为零。这些对应关系提示部分空白可能表示没有设施；缺少上传副本的原始字段说明，项目没有据此改写原值。

房价 Id 为 524 的记录中，改建年份 2008 晚于出售年份 2007。程序标记并保留该记录，主模型未使用改建年份。四条居住面积超过 4,000 的记录也保留。相关原值和定位行号见 \code{tables/row_flags.csv}。
\source{\code{tables/input_integrity.csv}；\code{tables/column_quality.csv}；\code{tables/house_missing_semantics.csv}}

\code{adult.data} 有 32,561 条，\code{adult.test} 有 16,281 条。原训练文件的全字段重复额外记录为 24 条，按 14 个原始非目标字段比较时为 25 条。原表缺少可靠的个人标识，相同字段可能来自重复录入，也可能属于不同个人。记录全部保留，划分时按全部原始预测字段的 SHA-256 指纹分组，相同字段的记录放在同一侧。

原训练、测试文件共有 26 个非目标字段指纹。原始测试文件已参与内容核查，没有用于模型训练和性能计算。以后若使用它评价模型，应事先规定重合指纹的处理方式。当前均值、比例和模型拟合都未使用 \code{fnlwgt} 加权。数据记录的是 1994 年相关人口普查样本的收入类别\citep{adult}。

上传表有 27 个赛季，从 1996--97 到 2022--23；共有 2,551 个不同姓名字符串。下列四组姓名与赛季键各有两条记录，年龄或身高不同：Marcus Williams 的 2007--08、2008--09，Chris Johnson 的 2012--13，以及 Tony Mitchell 的 2013--14。缺少稳定球员 ID，八条记录均保留，未按姓名与赛季强行合并。姓名分组实验将相同姓名放在同组，没有完成真实身份消歧。

\code{ts_pct} 有 15 条大于 1。真实投篮效率的分母是加权出手次数，不是普通二项试验总次数\citep{nbaglossary}。项目没有把该字段截断到 0--1；由于缺少对应的逐条出手次数，未对这些高值重新计算。基准特征未使用该字段。

数学、阅读、写作均在 0--100 分内，没有空白或全字段重复，均值分别为 66.089、69.169 和 68.054。源站致谢链接指向 \code{generated_data/exams}，原核查记录中的许可为 Unknown；本项目未取得生成脚本或真实抽样证明。本文将其作为课程样本分析，并保留源站记录\citep{student}。

\subsection{软件参数与统计约定}
\begin{center}\small
\begin{tabularx}{\textwidth}{lY}
\toprule
项目 & 当前实现\\
\midrule
筛选与分组 & 最多 8 个 AND 条件、2 个分组字段；展示前若干行不改变计算样本。\\
统计 & 计数、均值、中位数、和、极值、标准差；两数值变量的 Pearson 相关；直方图、散点图、分组柱形图。\\
分母 & 分别记录组内行数 \code{n_rows} 和非缺失数 \code{n_valid}；全缺失组的和返回缺失。\\
训练规模 & 每次一个目标、最多 30 个特征；有效目标 100 至 100,000 条，划分后训练至少 60、验证至少 20 条。\\
预处理 & 数值中位数插补、训练均值和标准差缩放；分类单独缺失类、独热编码，最多 30 个编码类别并保留参考类。\\
未见类别 & 编为全零并提示；训练中全缺失的特征报错。\\
逻辑回归 & $C=1$，\code{lbfgs}，最多 1,500 次迭代，阈值 0.5，无额外概率校准。\\
森林 & 120 棵树，深度 12，叶节点最少 5 条，种子 42，单线程；置换重要性重复 3 次。\\
会话预测 & 最多保留四个表格模型；重启后重新拟合，不能继续使用已失效的模型编号。\\
\bottomrule
\end{tabularx}
\end{center}
达到条数上限时，程序报错，不自动抽取一部分样本；缺失目标只在本次拟合中排除，不改写原文件。没有实现的算法、模型未收敛、非法特征和样本不足都返回具体原因。线性模型不输出 p 值，森林不输出预测区间。

Pearson 相关使用两列同时非缺失的记录。令这些记录的集合为 $A$，其均值为 $\bar x_A,\bar y_A$，则
\[
r=\frac{\sum_{i\in A}(x_i-\bar x_A)(y_i-\bar y_A)}
{\sqrt{\sum_{i\in A}(x_i-\bar x_A)^2\sum_{i\in A}(y_i-\bar y_A)^2}}.
\]
分子是共同变化，分母用各自波动进行标准化。项目至少要求三条成对有效记录；常数列无法得到有定义的相关系数。相关操作不进行因果识别或显著性检验。

\subsection{运行与检查记录}
四份表各有两个主要模型。房价和 NBA 各增加两个改变样本划分的模型，共 12 个固定配置。项目保存了每条验证记录的行号、实际值和预测值。原实验独立复算了 58 项指标。表~\ref{tab:houses} 第一行另取自两个已完成的演示模型，不计入这 12 个固定配置。

从保存的逐行预测可以复算连续目标的 MAE、RMSE 和 $R^2$，并从收入分类概率复算 AUC、$F_1$、Brier 分数和对数损失。每次都按行号核对预测记录与验证样本。复算直接使用已保存的预测。附带的 \field{check\_results.py} 读取原工程，输出相应数值。

\begin{center}\small
\begin{tabularx}{\textwidth}{lY}
\toprule
记录 & 实际完成内容\\
\midrule
0.2 数据核查 & 四份上传数据、缺失、重复、字段类型、文件一致性与划分问题。\\
0.2 固定模型 & 12 个配置；保存逐行验证预测；58 项指标独立复算一致。\\
0.2 组件检查 & 100 项通过，使用明确标注的合成夹具；模拟模型只检查消息协议。\\
0.2 实际数据接口 & 35 项通过，使用人工 JSON 参数；覆盖统计、训练、预测、PDF 及结果下载。\\
0.2 检索检查 & 16 道固定题前六条命中，MRR 为 0.96875；多数包含原字段名。\\
0.3 讲解计算 & 实算 CollgCr 均值和三条记录的 OLS；重建 292 条房价预测、6,512 条 Adult 概率；检查梯度与 Hessian。\\
仍需目标电脑实测 & DeepSeek API 方案生成质量、浏览器完整交互、Mac 响应时间与峰值内存。\\
\bottomrule
\end{tabularx}
\end{center}
DeepSeek 联网适配版的检查记录另见 \code{codes/verification/deepseek/scope.txt}：135 项测试通过，其中 26 项为新增联网适配测试；测试使用占位密钥和模拟 API 回复，没有调用付费服务。这些检查验证请求组织、错误处理和本地计算流程，不能据此报告真实 API 的方案生成准确率。

旧环境的 HTTP 状态接口可达，Chromium 打开页面却返回 \code{ERR_BLOCKED_BY_ADMINISTRATOR}，因此浏览器交互未验收。旧版全部运行记录随包保留。0.3 版另外复跑了 100 项组件测试，全部通过；35 项实际数据接口检查仍引用旧版记录，十二个模型配置没有重新拟合。

十二个配置的预处理参数来自各自训练子集；独立复算先按稳定行号对齐实际值和预测值，再重算指标。开发时出现过预测顺序与划分表顺序不一致，已通过行号对齐修正复算程序，原预测值未改。失败尝试保留在开发记录中。

\begin{center}\small
\begin{tabularx}{\textwidth}{lY}
\toprule
路径 & 内容\\
\midrule
\code{tables/model_comparison.csv} & 十二个模型配置的全部指标。\\
\code{tables/metric_reconciliation.csv} & 从验证预测独立复算的 58 项结果。\\
\code{tables/teaching/} & 均值样本、三条记录的 OLS、系数乘特征的逐项贡献。\\
\code{artifacts/benchmarks/} & 旧版模型方案、划分、预测、系数、参数和执行代码。\\
\code{artifacts/teaching/} & 本次社区均值计算的完整输出。\\
\code{verification/editorial/} & 本次讲解计算、数学核对、文档编译与检查记录。\\
\bottomrule
\end{tabularx}
\end{center}
数学核对没有重新估计旧版基准模型。房价预测重建最大绝对差为 $5.82\times10^{-11}$ 美元，Adult 概率最大绝对差为 $3.33\times10^{-16}$；这些数值反映浮点运算顺序差异。数值差分检查使用四条确定性代数输入，梯度和 Hessian 的最大误差分别为 $1.08\times10^{-11}$ 与 $6.23\times10^{-12}$。

\subsection{在 Mac 上准备与运行}
工程内已有四份处理数据。当前演示入口为项目根目录的 \code{TableLab-DeepSeek.command}。双击该文件，或在终端进入项目目录运行：
\begin{lstlisting}
bash TableLab-DeepSeek.command
\end{lstlisting}
启动器准备 Python 环境并启动本机网页。首次按终端提示隐藏输入 DeepSeek API Key，密钥默认存入当前用户的 Mac 钥匙串。浏览器打开 \url{http://127.0.0.1:8765}，演示结束在服务终端按 Control-C。

当前 API 客户端调用 \url{https://api.deepseek.com/chat/completions}，模型参数以 \code{tablelab/cloud/deepseek_transport.py} 中的配置为准。请求发送问题、字段名与类型、字段和方法说明及必要的上一轮方案；不整表上传原始 CSV。用户在问题或筛选条件中输入的内容也会随请求发送。返回值通过本地格式和字段检查后，交由用户核对，再由本机执行器计算。统计、模型训练与结果文件均在本机完成。生成新方案需要联网，API 用量由服务提供方计费。

打开网页，选「房价」，输入「CollgCr 社区房屋的平均成交价是多少？」。点击「检索并生成方案」，核对社区条件与统计方法，查看「检索依据」里的字段说明。点击「校验并执行」，对照 150 条有效记录和 197,965.77 美元的均价，再打开本轮的 CSV 与执行器源代码。

清空上下文后，要求用全部房屋记录、居住面积和整体质量做线性回归。确认方案里没有社区筛选，运行后让听众改一个条件。只改算法时，两次训练应使用同一批房屋。分别查看运行编号与验证误差。

远程汇报时共享本机浏览器窗口并保持网络连接。服务不可用时查看网页错误提示；随包的 HTML 回放只能查看保存的例题，不能生成新方案。

\subsection{重建文档}
图形使用本目录的 PDF。论文使用 XeLaTeX、BibTeX 和 natbib，在本目录运行 \code{bash build_paper.sh}，生成包含附录的 \code{paper.pdf}。

